\begingroup
\makeatletter\tagsleft@true\let\veqno\@@leqno\makeatother
\ArithmeticHead{342}{L}{ANALYSIS.\quad NACHLASS.}
\NumberedItem{3}
\par\indent Aus der Verbindung der Gleichungen (1) und (27) folgt nun leicht, wenn man die zweite derivirte Function durch $F''(\alpha,\beta,\gamma)$ bezeichnet
\begin{equation}
\alpha\beta F(\alpha,\beta,\gamma)-(\gamma-(\alpha+\beta+1)x)F'(\alpha,\beta,\gamma)-(x-xx)F''(\alpha,\beta,\gamma)=0.\tag*{(30)}
\end{equation}
\par\indent Es sei $P$ der Werth von $F(\alpha,\beta,\gamma)$ und $Q$ der Werth derselben Function, wenn man für $x$, $1-x$ setzt. Man hat dann:
\[\begin{aligned}
\alpha\beta P-(\gamma-(\alpha+\beta+1)x)\frac{dP}{dx}-(x-xx)\frac{d^2P}{dx^2}&=0,\\[.7em]
\alpha\beta Q-(\alpha+\beta+1-\gamma-(\alpha+\beta+1)x)\frac{dQ}{dx}-(x-xx)\frac{d^2Q}{dx^2}&=0,
\end{aligned}\]
\noindent also, für den Fall, wo $2\gamma=\alpha+\beta+1$ ist, welcher eine besondere Aufmerksamkeit verdient,
\par\smallskip\noindent{\small[S. 40]}\par
\[\begin{aligned}
\alpha\beta P-\gamma(1-2x)\frac{dP}{dx}-(x-xx)\frac{d^2P}{dx^2}&=0,\\[.7em]
\alpha\beta Q-\gamma(1-2x)\frac{dQ}{dx}-(x-xx)\frac{d^2Q}{dx^2}&=0.
\end{aligned}\]
\noindent Hieraus folgt
\clearpage\SourceLeaf{350}\ArithmeticHead{343}{R}{{\scriptsize ZUR THEORIE DER REIHE $F(\alpha,\beta,\gamma,x)$.\quad III.}}
\[-\frac{\gamma(1-2x)dx}{x-xx}=\frac{Qd^2P-Pd^2Q}{QdP-PdQ},\]
\noindent also
\begin{equation}
Q\frac{dP}{dx}-P\frac{dQ}{dx}=A(x-xx)^{-\gamma},\tag*{(31)}
\end{equation}
\noindent wo $A$ eine Constante ist, deren Werth weiterhin zu
\[\frac{\alpha{.}\beta{.}\alpha+1{.}\beta+1{.}\alpha+2{.}\beta+2}{\gamma{.}\gamma{.}\gamma+1{.}\gamma+1{.}\gamma+2}\mathpunct{\raisebox{.4ex}{.\,.\,.}}\]
\noindent bestimmt wird. Also z.\,B. für $\alpha=\frac12$, $\beta=\frac12$, $\gamma=1$
\[A=\frac14\cdot\frac98\cdot\frac{25}{24}\cdot\frac{49}{48}\mathpunct{\raisebox{.4ex}{.\,.\,.}}=\frac1\pi\mathpunct{\raisebox{.35ex}{.}}\]
\noindent Dieselbe kann auch allgemein durch
\[\begin{aligned}
&\frac{\alpha\beta}{\gamma}\Biggl(1+\frac{\gamma-\alpha{.}\gamma-\beta}{1{.}\gamma+1}\\[.7em]
&\qquad+\frac{\gamma-\alpha{.}\gamma+1-\alpha{.}\gamma-\beta{.}\gamma+1-\beta}{1{.}2{.}\gamma+1{.}\gamma+2}+\text{etc.}\Biggr)
\end{aligned}\]
\noindent vorgestellt werden.
\clearpage
\endgroup

